\section{Discussion}
This chapter discusses the empirical results presented above. The discussion is organized around the three questions that motivated the evaluation in Chapter~4: how useful perturbation-based methods are as practical alternatives to backpropagation, how WP compares with NP in practice, and whether IS-NP improves over existing NP variants. We first give a brief overview of the main findings for each of these questions. The following subsections then discuss each finding in more detail, before turning to the limitations of the thesis and directions for future work.

\subsection{Overview of the main findings}

The first main finding is that perturbation-based learning methods are capable of training nontrivial multilayer networks, but remain clearly less effective than backpropagation as task difficulty increases. On the sinus regression task, the perturbation methods follow the backpropagation baseline closely, and on California housing and MNIST they still learn meaningful predictors or representations. However, they generally converge more slowly and to worse final loss and score values than backpropagation. On CIFAR-10, the gap becomes much larger: the perturbation methods learn above chance level, but do not approach the performance of backpropagation. This directly addresses Objective~3 by showing both where the perturbation methods work in the present experiments and where they begin to break down.

The second main finding is that the results consistently favour node perturbation over weight perturbation. Across the experiments, WP has higher estimator variance, lower cosine similarity, and weaker learning performance than the node perturbation methods. This difference is visible on the simpler tasks and becomes clearer on the more difficult classification tasks. These results are consistent with the variance-reduction argument developed in Chapter~3, where node-level perturbations are shown to reduce estimator variance relative to weight-level perturbations under the matched perturbation construction. The experiments therefore provide empirical support for the theoretical comparison between WP and NP, which is part of Objective~2.

The third main finding is that IS-NP is the strongest perturbation method overall. IS-NP obtains the highest cosine similarity, the lowest estimator variance across all tasks, and the strongest results on the more difficult classification tasks. The scaled-input diagnostic also separates IS-NP from vanilla NP and fan-in NP: when the input range is scaled, IS-NP remains stable while the diagnostics for the other NP variants deteriorate. This supports the motivation for input-scaled perturbations developed in Chapter~3, where IS-NP was introduced to match the pre-activation noise induced by WP. The empirical results therefore support the second part of Objective~2: developing and evaluating a theoretically motivated improvement to existing node perturbation methods.

\subsection{Practical capability and limits of perturbation-based learning}

Taken together, the results show that perturbation-based methods are capable of genuine learning in multilayer networks, but that their practical range is limited compared with backpropagation. The task progression is informative because it shows where this limitation appears. On sinus regression, the perturbation methods remain close to BP, showing that scalar-feedback updates can work well when the learning problem is simple. On California housing and MNIST, the methods still learn meaningful predictors or representations, but the gap to BP becomes clearer. On CIFAR-10, the limitation becomes decisive: the methods learn above chance, but do not solve the task well in the MLP setup used here. The overall pattern is therefore not that perturbation methods fail outright, but that their usefulness decreases as the learning problem demands more accurate gradient information.

This trend is best explained by the quality of the gradient estimates. Backpropagation computes the exact gradient, while perturbation methods estimate it from random perturbations and a scalar loss signal. As the task and model become more complex, the useful update direction becomes harder to recover from such limited information. This is reflected in the cosine similarity diagnostics: the best perturbation estimates are roughly around \(0.75\) on sinus regression, around \(0.3\) on California housing, around \(0.2\) on MNIST, and close to \(0.05\) on CIFAR-10. Part of this decline is expected in higher-dimensional parameter spaces, where noisy estimates are more likely to be nearly orthogonal to the true gradient. However, this is precisely why directional alignment is informative: it measures whether the perturbation update contains enough signal to point in a useful descent direction. The raw variance values are less straightforward to compare across tasks, since they depend on factors such as loss scale, gradient magnitude, architecture, parameter dimension, and perturbation scale. Variance is therefore most useful for comparing methods within the same task, while cosine similarity gives the clearest indication of how informative the perturbation updates are across tasks.

The practical interpretation is that perturbation-based learning has a meaningful but limited operating range in the present setup. The fact that the methods can learn MNIST-level tasks suggests that scalar-feedback learning is not restricted to trivial problems, and may be useful in settings where backpropagation is unavailable or biologically implausible. At the same time, the weak directional alignment on harder tasks shows what must improve for these methods to scale: the updates need to contain a stronger and more reliable component in the true descent direction. This is the main empirical insight for the third research question: the experiments identify both where perturbation methods remain useful and where they begin to break down relative to BP.

\subsection{Weight versus node perturbation}

The comparison between WP and NP gives one of the clearest links between the theoretical and empirical parts of the thesis. In Chapter~3, we showed that WP contains additional weight-level randomness compared with the matched node-level estimator IS-NP. The estimator diagnostics show this effect very clearly in practice. WP has substantially higher estimator variance than IS-NP across the experiments, often by orders of magnitude. The same pattern also appears for the other NP variants: even though the formal variance theorem was derived specifically for IS-NP, vanilla NP and fan-in-scaled NP also have much lower variance than WP in the empirical results. This suggests that the variance disadvantage of perturbing individual weights is not only a theoretical construction, but a strong practical effect in the networks studied here.

This variance difference also appears to matter for learning performance. On the simpler tasks, WP still learns reasonably well, which shows that its higher variance is not always enough to prevent useful learning. In favourable settings, even noisy weight-level perturbations can provide a usable update direction. However, as the tasks become more difficult, the disadvantage becomes more visible. WP generally has weaker directional alignment and worse learning performance than the NP methods, especially on the image-classification tasks. This indicates that the additional variance is not just a diagnostic artifact: it affects the optimization dynamics when the learning problem requires more accurate and reliable update directions.

Taken together, these results substantially clarify the WP--NP relationship identified as incomplete in Chapter~2. The thesis does not only show empirically that NP performs better than WP in the tested tasks; it also provides a theoretical explanation for why this should happen. WP and IS-NP estimate the same perturbation-smoothed objective, but WP does so with higher coordinate-wise variance. The empirical results then show the same variance effect directly, and also show that it translates into weaker learning performance as task complexity increases. For the feedforward nonlinear multilayer networks studied here, this strongly supports NP as the preferable perturbation family when both WP and NP are viable. WP remains conceptually important, and may still be relevant in settings where parameter-level perturbations are more natural. However, when both WP and NP are viable, the results suggest that node-level perturbation is the more promising starting point for further development.

\subsection{IS-NP versus other node perturbation methods}

The comparison between IS-NP and the existing NP variants provides empirical support for the method developed in Chapter~3. As in the WP--NP comparison above, this separation is not equally visible on the simplest tasks, because all perturbation methods are still able to produce useful updates when the learning problem is favourable. The clearest separation appears on the more challenging tasks. There, IS-NP clearly performs better than fan-in-scaled NP, which in turn performs better than vanilla NP, both in the estimator diagnostics, the final test loss, and convergence speed. This is important because IS-NP was not introduced as an ad hoc modification, but derived from the theoretical relationship between WP and NP. The empirical results therefore suggest that the theoretical development in Chapter~3 leads to a practically useful perturbation method.

The most likely explanation is the input-scaling argument developed in Chapter~3. Vanilla NP uses a fixed perturbation scale, while fan-in-scaled NP uses a layer-width correction based on an expected input norm. IS-NP instead scales the perturbation by the actual incoming activation norm for each layer and input. This predicts the ordering observed most clearly on the harder tasks: IS-NP has the strongest directional alignment, fan-in-scaled NP is next, and vanilla NP is weakest, while the variance results follow the opposite order. The scaled-input diagnostic gives more direct support for this explanation because it tests what happens when the activation scale is changed in a controlled way. When the sinus input range is changed from ([-1,1]) to ([-5,5]), the underlying regression problem is not fundamentally changed, but the scale of the incoming activations is. This makes the assumptions behind vanilla NP and fan-in-scaled NP less accurate: a fixed perturbation scale no longer matches the activation scale, and a fan-in correction cannot account for the actual norm of the input. As predicted, their cosine similarity deteriorates while their estimator variance increases. IS-NP remains much more stable, which is exactly what one would expect from a method that adapts to the actual activation norm rather than relying on a fixed or expected scale.

At the same time, this comparison should not be interpreted as a formal proof that IS-NP always dominates the other NP variants. Chapter~3 proves a variance result for IS-NP relative to WP, but the comparison with vanilla NP and fan-in-scaled NP is based on scaling arguments and empirical evidence. Other architectures, normalization schemes, perturbation scales, or tasks could change the relative performance of the NP methods. The appropriate conclusion is therefore that IS-NP is the most promising node perturbation method tested in this thesis. The results strongly suggest that its input-dependent scaling is useful, and that the new method derived from the WP--NP analysis improves over the NP variants currently considered in the perturbation-learning literature. However, this does not prove that the same ordering will hold generally in all settings.

\subsection{Limitations}

One limitation is that all experiments were performed using MLPs and a limited set of supervised learning problems. This was useful for isolating the behaviour of the learning rules in a controlled setting, but it also limits how far the empirical conclusions can be generalized. Modern neural networks often rely on architectural features such as convolutions, attention, residual connections, or recurrent structure, and these may interact differently with perturbation-based learning. Similarly, the four tasks used in this thesis were chosen to span increasing difficulty, but they do not cover the full range of problems on which biologically plausible learning rules could be relevant.

A second limitation is that biological plausibility is evaluated at an abstract algorithmic level. The thesis studies whether learning rules based on stochastic perturbations, local information, and scalar feedback can train standard artificial neural networks. However, artificial neural networks remain highly simplified models of biological neural systems. Real neural circuits are recurrent, dynamic, spiking, chemically modulated, and embedded in bodies and environments. The results therefore support perturbation-based learning as a biologically motivated principle, but they do not show that the specific algorithms studied here are directly implementable in the brain.

A third limitation is that the experiments use fixed algorithmic choices. For practical reasons, the thesis had to choose specific baselines, batch sizes, initialization schemes and update rules. These choices were kept relatively simple in order to make the comparison between methods clear. However, perturbation-based learning may be sensitive to such design choices, and other variants could perform differently. The results should therefore be interpreted as an evaluation of a controlled set of perturbation methods, not as an exhaustive search over all possible implementations.

\subsection{Future work}

A natural direction for future work is to test perturbation-based learning on more architectures and more tasks. In particular, applying IS-NP and the other perturbation methods to CNNs, residual networks, recurrent networks, or transformer-like architectures would give a better understanding of whether the results found here extend beyond MLPs. Testing on a wider range of datasets would also help clarify where perturbation-based methods break down and which task properties make them more or less effective.

Another important direction is to improve the optimization behaviour of perturbation methods. Modern backpropagation-based training benefits greatly from techniques such as momentum, adaptive learning rates, learning-rate schedules, normalization, and architectural choices such as skip connections. Perturbation methods may benefit from analogous improvements, such as better baselines, adaptive perturbation scales, decaying noise schedules, layerwise perturbations, or optimizer variants inspired by momentum or Adam. At the same time, such modifications should be evaluated not only by their empirical performance, but also by whether they preserve the biological motivation of the learning rule.

A third direction is to test perturbation-based learning in more biologically realistic settings. Since artificial neural networks are only simplified abstractions of biological neural systems, it would be valuable to study whether similar learning principles can be implemented in recurrent, spiking, neuromorphic, or biological-substrate-inspired architectures. Systems such as CL1 illustrate that increasingly biological forms of computation are becoming experimentally relevant. Even if the exact algorithms studied in this thesis cannot be applied directly to such systems, testing related perturbation-and-feedback principles in more biological architectures would be highly informative.

Finally, future work should compare perturbation-based learning more broadly against other biologically plausible alternatives to backpropagation. This thesis focuses on WP, NP, fan-in-scaled NP, and IS-NP in depth, but methods such as feedback alignment, predictive coding, local losses, and the forward-forward algorithm represent other possible routes away from standard backpropagation. A systematic comparison across such methods would make it clearer how perturbation learning performs both as a practical training method and as a biologically plausible learning principle.